\begin{figure}[t]
  \centering
  \begin{minipage}[t]{0.98\columnwidth}
    \centering
    \includegraphics[width=\textwidth]{../results/analysis-lda/temporal_change24_all-100-700ms/figures/temporal_model_fits_binned-temporal_change24_all-100-700ms.png}
    \vspace{-2mm}
    \subcaption{All pairs (including numerosity 1)}
  \end{minipage}
  \vspace{1mm}
  \begin{minipage}[t]{0.98\columnwidth}
    \centering
    \includegraphics[width=\textwidth]{../results/analysis-lda/temporal_change24_all-100-700ms/figures-no1/temporal_model_fits_binned-temporal_change24_all-100-700ms-no1.png}
    \vspace{-2mm}
    \subcaption{Excluding pairs containing numerosity 1}
  \end{minipage}
  \caption{Time-binned model-brain correlations for the 24-condition
  prime-target analysis. Bars show mean Spearman correlation between
  neural and model RDMs within 100\,ms bins ($N = 24$; error bars:
  95\% CI). (a)~With all pairs, ANS and Target PI(1-4) show similar
  correlations, particularly in the 100-200\,ms bin. (b)~After
  excluding pairs containing numerosity~1, Target PI(2-4) exceeds ANS
  in early and mid-latency bins while ANS drops to near zero from
  200\,ms onward. Model-difference tests show Target PI(2-4) $>$ ANS
  (170-310\,ms, $p = 0.0006$; two-sided). The shift between panels
  suggests that the presence of numerosity~1 can disproportionately
  boost ANS model fits in this analysis.}
  \label{fig:temporal-model-comparison}
\end{figure}
